\documentclass[12pt]{article}

\usepackage{amsmath,amssymb,booktabs,longtable,array,hyperref}
\topmargin -.5cm
\textheight 21cm
\oddsidemargin -.125cm
\textwidth 17cm

\newcommand{\SO}{\mathrm{SO}}
\newcommand{\D}{\Delta}
\newcommand{\la}{\langle}
\newcommand{\ra}{\rangle}
\newcommand{\half}{\frac{1}{2}}

\begin{document}
\sloppy

\begin{center}
{\Large\bf Unique Exact Labels in Konishi Table 1}
\end{center}

\section{Goal and Counting Convention}

Take Table~1 as reproduced in
\texttt{notes/konishi\_linearized\_ansatz/konishi\_linearized\_ansatz.tex}
and parse it at the level of \emph{exact displayed labels}
\begin{equation}
[a,b,c]_{(s_L,s_R)}.
\end{equation}
If a term appears with an explicit coefficient, such as
\begin{equation}
2(s_L,s_R),
\end{equation}
that contributes multiplicity two to the corresponding exact label.

So for each displayed pair $\big([a,b,c],(s_L,s_R)\big)$ I define the total Table-1 multiplicity
\begin{equation}
N\big([a,b,c]_{(s_L,s_R)}\big)
=\sum_{\D_0} N_{\D_0}\big([a,b,c]_{(s_L,s_R)}\big).
\end{equation}
The phrase ``appears only once in the entire table'' will mean
\begin{equation}
N\big([a,b,c]_{(s_L,s_R)}\big)=1.
\end{equation}

This is slightly stronger than the weaker statement ``supported at a single weight''.
I record below the four labels that are supported at only one weight but with coefficient two.

The branching convention is the same bookkeeping convention used in
\texttt{notes/level1\_primary\_comparison/level1\_primary\_comparison.tex}:
\begin{itemize}
\item keep the compact $\SO(4)\subset \SO(1,4)$ spin label $(s_L,s_R)$ exactly as displayed in Table~1;
\item branch only the $\SO(6)$ irrep to $\SO(5)$.
\end{itemize}
So the ``$\SO(1,4)\times \SO(5)$ decomposition'' below should be read as the compact-label bookkeeping
\begin{equation}
[a,b,c]_{(s_L,s_R)}
\rightsquigarrow
\bigoplus_{\mu\in [a,b,c]\downarrow \SO(5)} (s_L,s_R;\mu),
\end{equation}
not as a claim that Table~1 by itself determines the full noncompact highest weight of the
$\SO(1,4)$ representation.

\paragraph{Quadratic Casimir conventions.}
For an $\SO(2,4)\times \SO(6)$ highest-weight label
\begin{equation}
[j_1,j_2,j_3]_{(s_1,s_2)}
\qquad\text{with conformal weight}\qquad \D,
\end{equation}
the quadratic Casimir is
\begin{equation}
C_2^{\SO(2,4)\times \SO(6)}
=
C_2^{\SO(2,4)}+C_2^{\SO(6)},
\end{equation}
with
\begin{equation}
C_2^{\SO(2,4)}
=
\D(\D-4)+2s_1(s_1+1)+2s_2(s_2+1),
\end{equation}
and
\begin{equation}
C_2^{\SO(6)}
=
\frac18\Big(
3j_1^2+4j_2^2+3j_3^2
+4j_1j_2+2j_1j_3+4j_2j_3
+12j_1+16j_2+12j_3
\Big).
\end{equation}

\section{Summary}

The exact scan gives:
\begin{itemize}
\item $201$ distinct exact labels $[a,b,c]_{(s_L,s_R)}$ in Table~1.
\item $88$ of them have total multiplicity one.
\item These $88$ labels occur only at
\begin{equation}
\D_0=\frac{9}{2},\ 5,\ \frac{11}{2},\ 6,\ \frac{13}{2},\ 7,\ \frac{15}{2}.
\end{equation}
\end{itemize}

The weight distribution is:
\begin{table}[h]
\centering
\begin{tabular}{cc}
\toprule
$\D_0$ & number of exact labels with total multiplicity $1$ \\
\midrule
$\frac{9}{2}$ & $4$ \\
$5$ & $8$ \\
$\frac{11}{2}$ & $20$ \\
$6$ & $24$ \\
$\frac{13}{2}$ & $20$ \\
$7$ & $8$ \\
$\frac{15}{2}$ & $4$ \\
\bottomrule
\end{tabular}
\caption{Weight distribution of the exact labels that appear only once in the full table.}
\label{tab:summary}
\end{table}

The motivating example
\begin{equation}
[2,0,2]_{(0,0)}
\end{equation}
is \emph{not} unique in this sense: it occurs once at $\D_0=4$ and once at $\D_0=8$, hence
\begin{equation}
N\big([2,0,2]_{(0,0)}\big)=2.
\end{equation}

\section{Unique Exact Labels and Their \texorpdfstring{$\SO(1,4)\times \SO(5)$}{SO(1,4)xSO(5)} Bookkeeping Decomposition}

\small
\begin{longtable}{c p{0.23\textwidth} >{\raggedright\arraybackslash}p{0.60\textwidth}}
\caption{All exact Table-1 labels with total multiplicity one, grouped by $\D_0$. The right column keeps the compact $\SO(4)$ spin label fixed and branches only the $\SO(6)$ irrep to $\SO(5)$.}
\label{tab:unique-exact-labels}\\
\toprule
$\D_0$ & Table 1 label & bookkeeping $\SO(1,4)\times \SO(5)$ decomposition \\
\midrule
\endfirsthead
\toprule
$\D_0$ & Table 1 label & bookkeeping $\SO(1,4)\times \SO(5)$ decomposition \\
\midrule
\endhead
\midrule
\endfoot
\bottomrule
\endlastfoot
$\frac{9}{2}$ & $[0,0,1]_{(2,\frac{1}{2})}$ & $((2,\frac{1}{2});\langle 1,0\rangle)$ \\
$\frac{9}{2}$ & $[0,0,3]_{(\frac{3}{2},0)}$ & $((\frac{3}{2},0);\langle 3,0\rangle)$ \\
$\frac{9}{2}$ & $[1,0,0]_{(\frac{1}{2},2)}$ & $((\frac{1}{2},2);\langle 1,0\rangle)$ \\
$\frac{9}{2}$ & $[3,0,0]_{(0,\frac{3}{2})}$ & $((0,\frac{3}{2});\langle 3,0\rangle)$ \\
$5$ & $[0,0,2]_{(1,0)}$ & $((1,0);\langle 2,0\rangle)$ \\
$5$ & $[0,0,2]_{(2,0)}$ & $((2,0);\langle 2,0\rangle)$ \\
$5$ & $[0,1,2]_{(\frac{3}{2},\frac{1}{2})}$ & $((\frac{3}{2},\frac{1}{2});\langle 2,0\rangle)\oplus ((\frac{3}{2},\frac{1}{2});\langle 2,1\rangle)$ \\
$5$ & $[1,0,3]_{(1,0)}$ & $((1,0);\langle 2,1\rangle)\oplus ((1,0);\langle 4,0\rangle)$ \\
$5$ & $[2,0,0]_{(0,1)}$ & $((0,1);\langle 2,0\rangle)$ \\
$5$ & $[2,0,0]_{(0,2)}$ & $((0,2);\langle 2,0\rangle)$ \\
$5$ & $[2,1,0]_{(\frac{1}{2},\frac{3}{2})}$ & $((\frac{1}{2},\frac{3}{2});\langle 2,0\rangle)\oplus ((\frac{1}{2},\frac{3}{2});\langle 2,1\rangle)$ \\
$5$ & $[3,0,1]_{(0,1)}$ & $((0,1);\langle 2,1\rangle)\oplus ((0,1);\langle 4,0\rangle)$ \\
$\frac{11}{2}$ & $[0,0,1]_{(\frac{3}{2},2)}$ & $((\frac{3}{2},2);\langle 1,0\rangle)$ \\
$\frac{11}{2}$ & $[0,0,3]_{(0,\frac{1}{2})}$ & $((0,\frac{1}{2});\langle 3,0\rangle)$ \\
$\frac{11}{2}$ & $[0,0,3]_{(1,\frac{1}{2})}$ & $((1,\frac{1}{2});\langle 3,0\rangle)$ \\
$\frac{11}{2}$ & $[0,1,1]_{(2,\frac{1}{2})}$ & $((2,\frac{1}{2});\langle 1,0\rangle)\oplus ((2,\frac{1}{2});\langle 1,1\rangle)$ \\
$\frac{11}{2}$ & $[0,1,3]_{(\frac{1}{2},0)}$ & $((\frac{1}{2},0);\langle 3,0\rangle)\oplus ((\frac{1}{2},0);\langle 3,1\rangle)$ \\
$\frac{11}{2}$ & $[0,2,1]_{(\frac{3}{2},0)}$ & $((\frac{3}{2},0);\langle 1,0\rangle)\oplus ((\frac{3}{2},0);\langle 1,1\rangle)\oplus ((\frac{3}{2},0);\langle 1,2\rangle)$ \\
$\frac{11}{2}$ & $[0,3,1]_{(0,\frac{1}{2})}$ & $((0,\frac{1}{2});\langle 1,0\rangle)\oplus ((0,\frac{1}{2});\langle 1,1\rangle)\oplus ((0,\frac{1}{2});\langle 1,2\rangle)\oplus ((0,\frac{1}{2});\langle 1,3\rangle)$ \\
$\frac{11}{2}$ & $[1,0,0]_{(2,\frac{3}{2})}$ & $((2,\frac{3}{2});\langle 1,0\rangle)$ \\
$\frac{11}{2}$ & $[1,0,2]_{(\frac{3}{2},0)}$ & $((\frac{3}{2},0);\langle 1,1\rangle)\oplus ((\frac{3}{2},0);\langle 3,0\rangle)$ \\
$\frac{11}{2}$ & $[1,0,2]_{(\frac{3}{2},1)}$ & $((\frac{3}{2},1);\langle 1,1\rangle)\oplus ((\frac{3}{2},1);\langle 3,0\rangle)$ \\
$\frac{11}{2}$ & $[1,1,0]_{(\frac{1}{2},2)}$ & $((\frac{1}{2},2);\langle 1,0\rangle)\oplus ((\frac{1}{2},2);\langle 1,1\rangle)$ \\
$\frac{11}{2}$ & $[1,1,2]_{(1,\frac{1}{2})}$ & $((1,\frac{1}{2});\langle 1,1\rangle)\oplus ((1,\frac{1}{2});\langle 3,0\rangle)\oplus ((1,\frac{1}{2});\langle 1,2\rangle)\oplus ((1,\frac{1}{2});\langle 3,1\rangle)$ \\
$\frac{11}{2}$ & $[1,2,0]_{(0,\frac{3}{2})}$ & $((0,\frac{3}{2});\langle 1,0\rangle)\oplus ((0,\frac{3}{2});\langle 1,1\rangle)\oplus ((0,\frac{3}{2});\langle 1,2\rangle)$ \\
$\frac{11}{2}$ & $[1,3,0]_{(\frac{1}{2},0)}$ & $((\frac{1}{2},0);\langle 1,0\rangle)\oplus ((\frac{1}{2},0);\langle 1,1\rangle)\oplus ((\frac{1}{2},0);\langle 1,2\rangle)\oplus ((\frac{1}{2},0);\langle 1,3\rangle)$ \\
$\frac{11}{2}$ & $[2,0,1]_{(0,\frac{3}{2})}$ & $((0,\frac{3}{2});\langle 1,1\rangle)\oplus ((0,\frac{3}{2});\langle 3,0\rangle)$ \\
$\frac{11}{2}$ & $[2,0,1]_{(1,\frac{3}{2})}$ & $((1,\frac{3}{2});\langle 1,1\rangle)\oplus ((1,\frac{3}{2});\langle 3,0\rangle)$ \\
$\frac{11}{2}$ & $[2,1,1]_{(\frac{1}{2},1)}$ & $((\frac{1}{2},1);\langle 1,1\rangle)\oplus ((\frac{1}{2},1);\langle 3,0\rangle)\oplus ((\frac{1}{2},1);\langle 1,2\rangle)\oplus ((\frac{1}{2},1);\langle 3,1\rangle)$ \\
$\frac{11}{2}$ & $[3,0,0]_{(\frac{1}{2},0)}$ & $((\frac{1}{2},0);\langle 3,0\rangle)$ \\
$\frac{11}{2}$ & $[3,0,0]_{(\frac{1}{2},1)}$ & $((\frac{1}{2},1);\langle 3,0\rangle)$ \\
$\frac{11}{2}$ & $[3,1,0]_{(0,\frac{1}{2})}$ & $((0,\frac{1}{2});\langle 3,0\rangle)\oplus ((0,\frac{1}{2});\langle 3,1\rangle)$ \\
$6$ & $[0,0,0]_{(2,2)}$ & $((2,2);\langle 0,0\rangle)$ \\
$6$ & $[0,0,2]_{(\frac{3}{2},\frac{3}{2})}$ & $((\frac{3}{2},\frac{3}{2});\langle 2,0\rangle)$ \\
$6$ & $[0,0,4]_{(0,0)}$ & $((0,0);\langle 4,0\rangle)$ \\
$6$ & $[0,1,2]_{(0,0)}$ & $((0,0);\langle 2,0\rangle)\oplus ((0,0);\langle 2,1\rangle)$ \\
$6$ & $[0,1,2]_{(1,1)}$ & $((1,1);\langle 2,0\rangle)\oplus ((1,1);\langle 2,1\rangle)$ \\
$6$ & $[0,2,0]_{(0,1)}$ & $((0,1);\langle 0,0\rangle)\oplus ((0,1);\langle 0,1\rangle)\oplus ((0,1);\langle 0,2\rangle)$ \\
$6$ & $[0,2,0]_{(0,2)}$ & $((0,2);\langle 0,0\rangle)\oplus ((0,2);\langle 0,1\rangle)\oplus ((0,2);\langle 0,2\rangle)$ \\
$6$ & $[0,2,0]_{(1,0)}$ & $((1,0);\langle 0,0\rangle)\oplus ((1,0);\langle 0,1\rangle)\oplus ((1,0);\langle 0,2\rangle)$ \\
$6$ & $[0,2,0]_{(2,0)}$ & $((2,0);\langle 0,0\rangle)\oplus ((2,0);\langle 0,1\rangle)\oplus ((2,0);\langle 0,2\rangle)$ \\
$6$ & $[0,2,2]_{(\frac{1}{2},\frac{1}{2})}$ & $((\frac{1}{2},\frac{1}{2});\langle 2,0\rangle)\oplus ((\frac{1}{2},\frac{1}{2});\langle 2,1\rangle)\oplus ((\frac{1}{2},\frac{1}{2});\langle 2,2\rangle)$ \\
$6$ & $[0,4,0]_{(0,0)}$ & $((0,0);\langle 0,0\rangle)\oplus ((0,0);\langle 0,1\rangle)\oplus ((0,0);\langle 0,2\rangle)\oplus ((0,0);\langle 0,3\rangle)\oplus ((0,0);\langle 0,4\rangle)$ \\
$6$ & $[1,0,1]_{(1,2)}$ & $((1,2);\langle 0,1\rangle)\oplus ((1,2);\langle 2,0\rangle)$ \\
$6$ & $[1,0,1]_{(2,1)}$ & $((2,1);\langle 0,1\rangle)\oplus ((2,1);\langle 2,0\rangle)$ \\
$6$ & $[1,0,3]_{(\frac{1}{2},\frac{1}{2})}$ & $((\frac{1}{2},\frac{1}{2});\langle 2,1\rangle)\oplus ((\frac{1}{2},\frac{1}{2});\langle 4,0\rangle)$ \\
$6$ & $[1,2,1]_{(0,0)}$ & $((0,0);\langle 0,1\rangle)\oplus ((0,0);\langle 2,0\rangle)\oplus ((0,0);\langle 0,2\rangle)\oplus ((0,0);\langle 2,1\rangle)\oplus ((0,0);\langle 0,3\rangle)\oplus ((0,0);\langle 2,2\rangle)$ \\
$6$ & $[1,2,1]_{(0,1)}$ & $((0,1);\langle 0,1\rangle)\oplus ((0,1);\langle 2,0\rangle)\oplus ((0,1);\langle 0,2\rangle)\oplus ((0,1);\langle 2,1\rangle)\oplus ((0,1);\langle 0,3\rangle)\oplus ((0,1);\langle 2,2\rangle)$ \\
$6$ & $[1,2,1]_{(1,0)}$ & $((1,0);\langle 0,1\rangle)\oplus ((1,0);\langle 2,0\rangle)\oplus ((1,0);\langle 0,2\rangle)\oplus ((1,0);\langle 2,1\rangle)\oplus ((1,0);\langle 0,3\rangle)\oplus ((1,0);\langle 2,2\rangle)$ \\
$6$ & $[2,0,0]_{(\frac{3}{2},\frac{3}{2})}$ & $((\frac{3}{2},\frac{3}{2});\langle 2,0\rangle)$ \\
$6$ & $[2,0,2]_{(1,1)}$ & $((1,1);\langle 0,2\rangle)\oplus ((1,1);\langle 2,1\rangle)\oplus ((1,1);\langle 4,0\rangle)$ \\
$6$ & $[2,1,0]_{(0,0)}$ & $((0,0);\langle 2,0\rangle)\oplus ((0,0);\langle 2,1\rangle)$ \\
$6$ & $[2,1,0]_{(1,1)}$ & $((1,1);\langle 2,0\rangle)\oplus ((1,1);\langle 2,1\rangle)$ \\
$6$ & $[2,2,0]_{(\frac{1}{2},\frac{1}{2})}$ & $((\frac{1}{2},\frac{1}{2});\langle 2,0\rangle)\oplus ((\frac{1}{2},\frac{1}{2});\langle 2,1\rangle)\oplus ((\frac{1}{2},\frac{1}{2});\langle 2,2\rangle)$ \\
$6$ & $[3,0,1]_{(\frac{1}{2},\frac{1}{2})}$ & $((\frac{1}{2},\frac{1}{2});\langle 2,1\rangle)\oplus ((\frac{1}{2},\frac{1}{2});\langle 4,0\rangle)$ \\
$6$ & $[4,0,0]_{(0,0)}$ & $((0,0);\langle 4,0\rangle)$ \\
$\frac{13}{2}$ & $[0,0,1]_{(2,\frac{3}{2})}$ & $((2,\frac{3}{2});\langle 1,0\rangle)$ \\
$\frac{13}{2}$ & $[0,0,3]_{(\frac{1}{2},0)}$ & $((\frac{1}{2},0);\langle 3,0\rangle)$ \\
$\frac{13}{2}$ & $[0,0,3]_{(\frac{1}{2},1)}$ & $((\frac{1}{2},1);\langle 3,0\rangle)$ \\
$\frac{13}{2}$ & $[0,1,1]_{(\frac{1}{2},2)}$ & $((\frac{1}{2},2);\langle 1,0\rangle)\oplus ((\frac{1}{2},2);\langle 1,1\rangle)$ \\
$\frac{13}{2}$ & $[0,1,3]_{(0,\frac{1}{2})}$ & $((0,\frac{1}{2});\langle 3,0\rangle)\oplus ((0,\frac{1}{2});\langle 3,1\rangle)$ \\
$\frac{13}{2}$ & $[0,2,1]_{(0,\frac{3}{2})}$ & $((0,\frac{3}{2});\langle 1,0\rangle)\oplus ((0,\frac{3}{2});\langle 1,1\rangle)\oplus ((0,\frac{3}{2});\langle 1,2\rangle)$ \\
$\frac{13}{2}$ & $[0,3,1]_{(\frac{1}{2},0)}$ & $((\frac{1}{2},0);\langle 1,0\rangle)\oplus ((\frac{1}{2},0);\langle 1,1\rangle)\oplus ((\frac{1}{2},0);\langle 1,2\rangle)\oplus ((\frac{1}{2},0);\langle 1,3\rangle)$ \\
$\frac{13}{2}$ & $[1,0,0]_{(\frac{3}{2},2)}$ & $((\frac{3}{2},2);\langle 1,0\rangle)$ \\
$\frac{13}{2}$ & $[1,0,2]_{(0,\frac{3}{2})}$ & $((0,\frac{3}{2});\langle 1,1\rangle)\oplus ((0,\frac{3}{2});\langle 3,0\rangle)$ \\
$\frac{13}{2}$ & $[1,0,2]_{(1,\frac{3}{2})}$ & $((1,\frac{3}{2});\langle 1,1\rangle)\oplus ((1,\frac{3}{2});\langle 3,0\rangle)$ \\
$\frac{13}{2}$ & $[1,1,0]_{(2,\frac{1}{2})}$ & $((2,\frac{1}{2});\langle 1,0\rangle)\oplus ((2,\frac{1}{2});\langle 1,1\rangle)$ \\
$\frac{13}{2}$ & $[1,1,2]_{(\frac{1}{2},1)}$ & $((\frac{1}{2},1);\langle 1,1\rangle)\oplus ((\frac{1}{2},1);\langle 3,0\rangle)\oplus ((\frac{1}{2},1);\langle 1,2\rangle)\oplus ((\frac{1}{2},1);\langle 3,1\rangle)$ \\
$\frac{13}{2}$ & $[1,2,0]_{(\frac{3}{2},0)}$ & $((\frac{3}{2},0);\langle 1,0\rangle)\oplus ((\frac{3}{2},0);\langle 1,1\rangle)\oplus ((\frac{3}{2},0);\langle 1,2\rangle)$ \\
$\frac{13}{2}$ & $[1,3,0]_{(0,\frac{1}{2})}$ & $((0,\frac{1}{2});\langle 1,0\rangle)\oplus ((0,\frac{1}{2});\langle 1,1\rangle)\oplus ((0,\frac{1}{2});\langle 1,2\rangle)\oplus ((0,\frac{1}{2});\langle 1,3\rangle)$ \\
$\frac{13}{2}$ & $[2,0,1]_{(\frac{3}{2},0)}$ & $((\frac{3}{2},0);\langle 1,1\rangle)\oplus ((\frac{3}{2},0);\langle 3,0\rangle)$ \\
$\frac{13}{2}$ & $[2,0,1]_{(\frac{3}{2},1)}$ & $((\frac{3}{2},1);\langle 1,1\rangle)\oplus ((\frac{3}{2},1);\langle 3,0\rangle)$ \\
$\frac{13}{2}$ & $[2,1,1]_{(1,\frac{1}{2})}$ & $((1,\frac{1}{2});\langle 1,1\rangle)\oplus ((1,\frac{1}{2});\langle 3,0\rangle)\oplus ((1,\frac{1}{2});\langle 1,2\rangle)\oplus ((1,\frac{1}{2});\langle 3,1\rangle)$ \\
$\frac{13}{2}$ & $[3,0,0]_{(0,\frac{1}{2})}$ & $((0,\frac{1}{2});\langle 3,0\rangle)$ \\
$\frac{13}{2}$ & $[3,0,0]_{(1,\frac{1}{2})}$ & $((1,\frac{1}{2});\langle 3,0\rangle)$ \\
$\frac{13}{2}$ & $[3,1,0]_{(\frac{1}{2},0)}$ & $((\frac{1}{2},0);\langle 3,0\rangle)\oplus ((\frac{1}{2},0);\langle 3,1\rangle)$ \\
$7$ & $[0,0,2]_{(0,1)}$ & $((0,1);\langle 2,0\rangle)$ \\
$7$ & $[0,0,2]_{(0,2)}$ & $((0,2);\langle 2,0\rangle)$ \\
$7$ & $[0,1,2]_{(\frac{1}{2},\frac{3}{2})}$ & $((\frac{1}{2},\frac{3}{2});\langle 2,0\rangle)\oplus ((\frac{1}{2},\frac{3}{2});\langle 2,1\rangle)$ \\
$7$ & $[1,0,3]_{(0,1)}$ & $((0,1);\langle 2,1\rangle)\oplus ((0,1);\langle 4,0\rangle)$ \\
$7$ & $[2,0,0]_{(1,0)}$ & $((1,0);\langle 2,0\rangle)$ \\
$7$ & $[2,0,0]_{(2,0)}$ & $((2,0);\langle 2,0\rangle)$ \\
$7$ & $[2,1,0]_{(\frac{3}{2},\frac{1}{2})}$ & $((\frac{3}{2},\frac{1}{2});\langle 2,0\rangle)\oplus ((\frac{3}{2},\frac{1}{2});\langle 2,1\rangle)$ \\
$7$ & $[3,0,1]_{(1,0)}$ & $((1,0);\langle 2,1\rangle)\oplus ((1,0);\langle 4,0\rangle)$ \\
$\frac{15}{2}$ & $[0,0,1]_{(\frac{1}{2},2)}$ & $((\frac{1}{2},2);\langle 1,0\rangle)$ \\
$\frac{15}{2}$ & $[0,0,3]_{(0,\frac{3}{2})}$ & $((0,\frac{3}{2});\langle 3,0\rangle)$ \\
$\frac{15}{2}$ & $[1,0,0]_{(2,\frac{1}{2})}$ & $((2,\frac{1}{2});\langle 1,0\rangle)$ \\
$\frac{15}{2}$ & $[3,0,0]_{(\frac{3}{2},0)}$ & $((\frac{3}{2},0);\langle 3,0\rangle)$ \\

\end{longtable}
\normalsize

\section{Single-Weight But Non-Unique Labels}

If one weakens the criterion and asks only for support at a single weight, four additional labels appear.
Each of them lives only at $\D_0=6$, but with coefficient two rather than one:
\begin{itemize}
\item $[0,1,0]_{(\frac{3}{2},\frac{3}{2})}$:
\begin{equation}
((\tfrac{3}{2},\tfrac{3}{2});\la 0,0\ra)\oplus ((\tfrac{3}{2},\tfrac{3}{2});\la 0,1\ra).
\end{equation}
\item $[0,3,0]_{(\frac{1}{2},\frac{1}{2})}$:
\begin{equation}
((\tfrac{1}{2},\tfrac{1}{2});\la 0,0\ra)\oplus ((\tfrac{1}{2},\tfrac{1}{2});\la 0,1\ra)\oplus
((\tfrac{1}{2},\tfrac{1}{2});\la 0,2\ra)\oplus ((\tfrac{1}{2},\tfrac{1}{2});\la 0,3\ra).
\end{equation}
\item $[1,1,1]_{(\frac{1}{2},\frac{3}{2})}$:
\begin{equation}
((\tfrac{1}{2},\tfrac{3}{2});\la 0,1\ra)\oplus ((\tfrac{1}{2},\tfrac{3}{2});\la 2,0\ra)\oplus
((\tfrac{1}{2},\tfrac{3}{2});\la 0,2\ra)\oplus ((\tfrac{1}{2},\tfrac{3}{2});\la 2,1\ra).
\end{equation}
\item $[1,1,1]_{(\frac{3}{2},\frac{1}{2})}$:
\begin{equation}
((\tfrac{3}{2},\tfrac{1}{2});\la 0,1\ra)\oplus ((\tfrac{3}{2},\tfrac{1}{2});\la 2,0\ra)\oplus
((\tfrac{3}{2},\tfrac{1}{2});\la 0,2\ra)\oplus ((\tfrac{3}{2},\tfrac{1}{2});\la 2,1\ra).
\end{equation}
\end{itemize}

\section{Bookkeeping-Isolated \texorpdfstring{$\SO(1,4)\times \SO(5)$}{SO(1,4)xSO(5)} Labels}

One can ask a stronger question than simple exact-label uniqueness. For each exact Table-1 label
\begin{equation}
[a,b,c]_{(s_L,s_R)},
\end{equation}
branch it to the bookkeeping set of labels
\begin{equation}
(s_L,s_R;\mu),\qquad \mu\in [a,b,c]\downarrow \SO(5),
\end{equation}
and ask whether any of these bookkeeping $\SO(1,4)\times \SO(5)$ labels are produced by
some \emph{other} exact Table-1 label.

With this convention there is exactly one \emph{fully isolated} exact label:
\begin{equation}
[0,0,0]_{(2,2)}\qquad \text{at}\qquad \D_0=6,
\end{equation}
because it branches only to
\begin{equation}
((2,2);\la 0,0\ra),
\end{equation}
and no other exact Table-1 label contributes this bookkeeping rep.

There are also ten \emph{partially isolated} exact labels: they do share some bookkeeping reps
with other labels, but each of them has at least one bookkeeping component that is unique in the
full Konishi table. They are:
\begin{itemize}
\item $[0,1,0]_{(1,2)}$ at $\D_0=5,7$, with unique component
\begin{equation}
((1,2);\la 0,0\ra).
\end{equation}
\item $[0,1,0]_{(2,1)}$ at $\D_0=5,7$, with unique component
\begin{equation}
((2,1);\la 0,0\ra).
\end{equation}
\item $[0,2,0]_{(0,2)}$ at $\D_0=6$, with unique components
\begin{equation}
((0,2);\la 0,1\ra),\qquad ((0,2);\la 0,2\ra).
\end{equation}
\item $[0,2,0]_{(2,0)}$ at $\D_0=6$, with unique components
\begin{equation}
((2,0);\la 0,1\ra),\qquad ((2,0);\la 0,2\ra).
\end{equation}
\item $[0,4,0]_{(0,0)}$ at $\D_0=6$, with unique component
\begin{equation}
((0,0);\la 0,4\ra).
\end{equation}
\item $[1,0,1]_{(1,2)}$ at $\D_0=6$, with unique component
\begin{equation}
((1,2);\la 2,0\ra).
\end{equation}
\item $[1,0,1]_{(2,1)}$ at $\D_0=6$, with unique component
\begin{equation}
((2,1);\la 2,0\ra).
\end{equation}
\item $[1,2,1]_{(0,1)}$ at $\D_0=6$, with unique component
\begin{equation}
((0,1);\la 2,2\ra).
\end{equation}
\item $[1,2,1]_{(1,0)}$ at $\D_0=6$, with unique component
\begin{equation}
((1,0);\la 2,2\ra).
\end{equation}
\item $[2,0,2]_{(1,1)}$ at $\D_0=6$, with unique component
\begin{equation}
((1,1);\la 4,0\ra).
\end{equation}
\end{itemize}

So the answer to the stronger question is:
\begin{itemize}
\item exactly one exact Table-1 label is completely isolated after branching,
namely $[0,0,0]_{(2,2)}$;
\item ten further exact labels are only partially isolated, in the sense that they each carry
at least one bookkeeping $\SO(1,4)\times \SO(5)$ label that no other exact Table-1 label carries.
\end{itemize}

\section{NSNS Vertex Operator for the Fully Isolated Mode}

For the fully isolated mode
\begin{equation}
[0,0,0]_{(2,2)}\qquad \text{at}\qquad \D_0=6,
\end{equation}
the natural highest-weight NSNS representative is the same spin-4 singlet seed used in
\texttt{notes/konishi\_linearized\_ansatz/konishi\_linearized\_ansatz.tex}.
Introduce
\begin{equation}
X^\pm=X^1\pm iX^2,
\end{equation}
and write the massive NSNS vertex in the form
\begin{equation}
V^{(-1,-1)}=
H_{\mu\nu;\rho\sigma}(X)\,
c\bar c\,
e^{-\phi}\psi^\mu\partial X^\nu\,
e^{-\bar\phi}\bar\psi^\rho\bar\partial X^\sigma,
\end{equation}
with $H_{\mu\nu;\rho\sigma}$ in the symmetric-traceless spin-4 representation, i.e.
\begin{equation}
H_{\mu\nu;\rho\sigma}
=H_{(\mu\nu);(\rho\sigma)},\qquad
\eta^{\mu\nu}H_{\mu\nu;\rho\sigma}=0,\qquad
\eta^{\rho\sigma}H_{\mu\nu;\rho\sigma}=0,
\end{equation}
and with the left and right index pairs projected to the highest-weight $(2,2)$ component.

An explicit highest-weight tensor structure is obtained from a null polarization vector in the chosen
$AdS_5$ complex plane,
\begin{equation}
n^\mu=(1,i,0,0,0),\qquad n_\mu n^\mu=0,
\end{equation}
for which one may write
\begin{equation}
H_{\mu\nu;\rho\sigma}(X)=h(X)\,n_\mu n_\nu n_\rho n_\sigma.
\end{equation}
Because $n^2=0$, this tensor is automatically symmetric and traceless in each index pair, so it is
the null-polarization representative of the spin-4 highest weight.

The explicit highest-weight representative is therefore
\begin{equation}
V^{(-1,-1)}_{[0,0,0]_{(2,2)}}
=H_{++;++}(X)\,
c\bar c\,
e^{-\phi}\psi^+\partial X^+\,
e^{-\bar\phi}\bar\psi^+\bar\partial X^+.
\end{equation}

Here the spin quantum number is carried by the polarization tensor component $H_{++;++}$:
the coefficient function is \emph{not} itself a spin-4 harmonic. Equivalently, $H_{++;++}(X)$
is just one highest-weight component of the full spin-4 field $H_{\mu\nu;\rho\sigma}(X)$.
With the null-polarization choice above, this component is simply
\begin{equation}
H_{++;++}(X)=h(X).
\end{equation}
No explicit coordinate ansatz for $h(X)$ is needed here. In particular, one should \emph{not}
identify $h(X)$ with a non-scalar monomial such as $(X^+)^4$ when the goal is only to specify the
representation content of the vertex. The point of the present notation is just that $h(X)$ is the
coefficient function multiplying the chosen highest-weight tensor component of the full spin-$4$ field.

\paragraph{Why this final vertex is spin $4$.}
The spin is not inferred from any particular coordinate ansatz for $h(X)$. Rather, the relevant
$AdS_5$ generators split as
\begin{equation}
J_{AB}=L_{AB}+\Sigma_{AB},
\end{equation}
where $L_{AB}$ acts on the coordinate dependence $H(X)$ and $\Sigma_{AB}$ acts on the tensor
indices carried by the NSNS oscillator structure. The local worldsheet factor
\begin{equation}
\psi^+\partial X^+\;\bar\psi^+\bar\partial X^+
\end{equation}
is precisely the highest-weight component of the closed-string NSNS tensor polarization.
This does \emph{not} mean that the whole state ``spins only in one plane.'' The label
\begin{equation}
(2,2)
\end{equation}
is an irrep of
\begin{equation}
\SO(4)\simeq SU(2)_L\times SU(2)_R,
\end{equation}
so the full spin-$4$ field contains all components generated from the highest-weight vector by the
lowering operators of $SU(2)_L\times SU(2)_R$. Writing only the $++++$ component is therefore a
basis choice for the highest-weight representative, not a statement that the full irrep is supported
only in the single complex plane used to define $X^\pm$.

It is also important that the symbol `$+$' here is \emph{not} ``one of the two spins.'' Rather, it is a
highest-weight vector of the $SO(4)$ vector representation itself. Under
\begin{equation}
\mathbf 4 \cong \left(\tfrac12,\tfrac12\right),
\end{equation}
a single vector index already carries weight in both $SU(2)$ factors at once. Concretely, if one chooses
\begin{equation}
J_L^3=\frac12(M_{12}+M_{34}),\qquad
J_R^3=\frac12(M_{12}-M_{34}),
\end{equation}
and takes
\begin{equation}
e_+=e_1+i e_2,
\end{equation}
then
\begin{equation}
M_{12}e_+=e_+,\qquad M_{34}e_+=0,
\end{equation}
so
\begin{equation}
J_L^3 e_+=\frac12 e_+,\qquad
J_R^3 e_+=\frac12 e_+.
\end{equation}
Thus one `$+$' index already has weight
\begin{equation}
\left(\tfrac12,\tfrac12\right).
\end{equation}
Therefore
\begin{equation}
\psi^+\partial X^+
\end{equation}
has highest weight
\begin{equation}
(1,1),
\end{equation}
and the closed-string product
\begin{equation}
\psi^+\partial X^+\;\bar\psi^+\bar\partial X^+
\end{equation}
is the highest-weight component of the spin label
\begin{equation}
(2,2).
\end{equation}
If one wanted to make this completely manifest, one should really use $SO(4)$ bispinor notation
$X^{\alpha\dot\alpha}$ and write the highest-weight vector as $X^{+\dot +}$; the present `$+$' notation is
just a compressed shorthand for that choice.

In the flat-space first-massive decomposition, the relevant left-moving NS representation is
\begin{equation}
\mathbf{44}\supset (1,1;\la 0,0\ra),
\end{equation}
while the three-form representation $\mathbf{84}$ contains no $(1,1;\la 0,0\ra)$ piece. Thus the
closed-string spin-$4$ singlet comes uniquely from
\begin{equation}
(1,1;\la 0,0\ra)_L\otimes (1,1;\la 0,0\ra)_R
\subset \mathbf{44}_L\otimes \mathbf{44}_R,
\end{equation}
which is exactly the unique flat-space NSNS contribution to
\begin{equation}
((2,2);\la 0,0\ra).
\end{equation}
So the vertex is spin $4$ because its oscillator/tensor structure lies in the
$\mathbf{44}_L\otimes \mathbf{44}_R$ spin-$(2,2)$ channel. The spin assignment is therefore a statement
about the tensor/oscillator representation content of the vertex, not about any particular non-scalar
coordinate ansatz for the coefficient function $h(X)$.

This is in the NSNS sector, carries the highest-weight $(2,2)$ polarization in the chosen
$AdS_5$ complex plane, and has trivial sphere profile
\begin{equation}
Y_{(0)}=1,
\end{equation}
so its bookkeeping $\SO(1,4)\times \SO(5)$ label is precisely
\begin{equation}
((2,2);\la 0,0\ra).
\end{equation}

In other words, among the isolated modes found above, the completely isolated one is represented
at leading order by a pure-$AdS$ NSNS spin-4 vertex with no extra $S^5$ harmonic turned on.
The profile is scalar under $\SO(5)$; the $(2,2)$ label comes from the $AdS$ highest-weight tensor
structure.

The scan underlying this note uses the same parser and $\SO(6)\to \SO(5)$ branching code as
\texttt{scripts/primary\_branching\_scan.py}.

\end{document}
